% ---------------------------------------------------------------------------
% Preamble for the CC5513 lecture notes
% ---------------------------------------------------------------------------
\usepackage[T1]{fontenc}
\usepackage[utf8]{inputenc}
\usepackage{lmodern}
\usepackage[a4paper,textwidth=155mm,top=27mm,bottom=30mm,hmarginratio=1:1,headsep=9mm]{geometry}
\usepackage{microtype}
\usepackage{amsmath,amssymb,amsthm,mathtools,bm}
\usepackage{graphicx}
\usepackage{booktabs,tabularx,multirow,array}
\usepackage[dvipsnames]{xcolor}
\usepackage{tikz}
\usetikzlibrary{arrows.meta,calc,positioning,decorations.markings,shapes.geometric,fit,backgrounds}
\usepackage[ruled,vlined,linesnumbered]{algorithm2e}
\usepackage{listings}
\usepackage[most]{tcolorbox}
\usepackage{enumitem}
\usepackage[font=small,labelfont=bf,margin=4mm]{caption}
\usepackage{subcaption}
\usepackage[numbers,sort&compress]{natbib}
\usepackage{fancyhdr}
\usepackage[colorlinks=true,linkcolor=MidnightBlue,citecolor=OliveGreen,urlcolor=RoyalBlue]{hyperref}
\usepackage[capitalise,noabbrev]{cleveref}

\graphicspath{{figs/}}
\setlist{itemsep=2pt,topsep=4pt}
\setlength{\parskip}{2pt}
\renewcommand{\bibname}{References}

% ------------------------------------------------------------------ headers
\pagestyle{fancy}
\fancyhf{}
\fancyhead[LE]{\small\nouppercase{\leftmark}}
\fancyhead[RO]{\small\nouppercase{\rightmark}}
\fancyfoot[C]{\small\thepage}
\renewcommand{\headrulewidth}{0.3pt}
\fancypagestyle{plain}{\fancyhf{}\fancyfoot[C]{\small\thepage}\renewcommand{\headrulewidth}{0pt}}

% ------------------------------------------------------------------ theorems
\theoremstyle{plain}
\newtheorem{theorem}{Theorem}[chapter]
\newtheorem{proposition}[theorem]{Proposition}
\newtheorem{lemma}[theorem]{Lemma}
\newtheorem{corollary}[theorem]{Corollary}
\theoremstyle{definition}
\newtheorem{definition}[theorem]{Definition}
\newtheorem{exercise}{Exercise}[chapter]
\theoremstyle{remark}
\newtheorem{remark}[theorem]{Remark}

\tcbset{colback=white,boxrule=0.5pt,arc=1.5pt,left=4pt,right=4pt,top=3pt,bottom=3pt,
        fonttitle=\bfseries\small,breakable}
\newtcbtheorem[number within=chapter]{example}{Example}%
  {colframe=MidnightBlue!70,colbacktitle=MidnightBlue!8,coltitle=black,colback=MidnightBlue!2}{ex}
\newtcolorbox{notebook}[1]{colframe=OliveGreen!70,colback=OliveGreen!3,colbacktitle=OliveGreen!10,
  coltitle=black,title={Companion notebook: \texttt{#1}}}
\newtcolorbox{keypoints}{colframe=BrickRed!60,colback=BrickRed!2,colbacktitle=BrickRed!8,
  coltitle=black,title={Key points}}

% ------------------------------------------------------------------ algorithms
\SetAlgoSkip{medskip}
\SetKwInOut{Input}{Input}
\SetKwInOut{Output}{Output}
\SetAlFnt{\small}
\SetAlCapFnt{\small}
\crefname{algocf}{Algorithm}{Algorithms}
\Crefname{algocf}{Algorithm}{Algorithms}
\crefname{example}{Example}{Examples}
\crefname{tcb@cnt@example}{Example}{Examples}
\crefname{exercise}{Exercise}{Exercises}

% ------------------------------------------------------------------ listings
\lstset{basicstyle=\ttfamily\footnotesize,keywordstyle=\color{MidnightBlue}\bfseries,
  commentstyle=\color{OliveGreen},stringstyle=\color{BrickRed},showstringspaces=false,
  frame=single,rulecolor=\color{black!25},breaklines=true,columns=fullflexible,
  language=Python,xleftmargin=2pt,xrightmargin=2pt,aboveskip=6pt,belowskip=4pt}

% ------------------------------------------------------------------ notation
\newcommand{\R}{\mathbb{R}}
\newcommand{\N}{\mathbb{N}}
\newcommand{\C}{\mathbb{C}}
\newcommand{\M}{\mathcal{M}}
\newcommand{\NN}{\mathcal{N}}
\newcommand{\bx}{\mathbf{x}}
\newcommand{\by}{\mathbf{y}}
\newcommand{\bz}{\mathbf{z}}
\newcommand{\bp}{\mathbf{p}}
\newcommand{\bq}{\mathbf{q}}
\newcommand{\bn}{\mathbf{n}}
\newcommand{\bu}{\mathbf{u}}
\newcommand{\bv}{\mathbf{v}}
\newcommand{\bc}{\mathbf{c}}
\newcommand{\bd}{\mathbf{d}}
\newcommand{\be}{\mathbf{e}}
\newcommand{\bt}{\mathbf{t}}
\newcommand{\bb}{\mathbf{b}}
\newcommand{\ba}{\mathbf{a}}
\newcommand{\br}{\mathbf{r}}
\newcommand{\bs}{\mathbf{s}}
\newcommand{\bo}{\mathbf{o}}
\newcommand{\bw}{\mathbf{w}}
\newcommand{\bfun}{\mathbf{f}}
\newcommand{\bg}{\mathbf{g}}
\newcommand{\bzero}{\mathbf{0}}
\newcommand{\mA}{\mathbf{A}}
\newcommand{\mB}{\mathbf{B}}
\newcommand{\mC}{\mathbf{C}}
\newcommand{\mD}{\mathbf{D}}
\newcommand{\mE}{\mathbf{E}}
\newcommand{\mG}{\mathbf{G}}
\newcommand{\mH}{\mathbf{H}}
\newcommand{\mI}{\mathbf{I}}
\newcommand{\mJ}{\mathbf{J}}
\newcommand{\mK}{\mathbf{K}}
\newcommand{\mL}{\mathbf{L}}
\newcommand{\mM}{\mathbf{M}}
\newcommand{\mP}{\mathbf{P}}
\newcommand{\mQ}{\mathbf{Q}}
\newcommand{\mR}{\mathbf{R}}
\newcommand{\mS}{\mathbf{S}}
\newcommand{\mT}{\mathbf{T}}
\newcommand{\mU}{\mathbf{U}}
\newcommand{\mV}{\mathbf{V}}
\newcommand{\mW}{\mathbf{W}}
\newcommand{\mX}{\mathbf{X}}
\newcommand{\mY}{\mathbf{Y}}
\newcommand{\mPhi}{\boldsymbol{\Phi}}
\newcommand{\mPsi}{\boldsymbol{\Psi}}
\newcommand{\mLambda}{\boldsymbol{\Lambda}}
\newcommand{\mPi}{\boldsymbol{\Pi}}
\newcommand{\mSigma}{\boldsymbol{\Sigma}}
\newcommand{\bdelta}{\boldsymbol{\delta}}
\newcommand{\bphi}{\boldsymbol{\phi}}
\newcommand{\Lap}{\Delta}
\newcommand{\T}{^{\mathsf{T}}}
\newcommand{\norm}[1]{\left\lVert #1\right\rVert}
\newcommand{\abs}[1]{\left\lvert #1\right\rvert}
\newcommand{\inner}[2]{\left\langle #1,#2\right\rangle}
\newcommand{\dd}{\,\mathrm{d}}
\DeclareMathOperator*{\argmin}{arg\,min}
\DeclareMathOperator*{\argmax}{arg\,max}
\DeclareMathOperator{\tr}{tr}
\DeclareMathOperator{\diag}{diag}
\DeclareMathOperator{\Div}{div}
\DeclareMathOperator{\area}{area}
\DeclareMathOperator{\vol}{vol}
\DeclareMathOperator{\sgn}{sgn}
\DeclareMathOperator{\rank}{rank}
\DeclareMathOperator{\atanh}{atanh}
\newcommand{\grad}{\nabla}

% figure helper: \fig[width]{file}{caption}{label}
\newcommand{\fig}[4][\textwidth]{%
  \begin{figure}[tbp]\centering\includegraphics[width=#1]{#2}\caption{#3}\label{#4}\end{figure}}
% figure at natural size
\newcommand{\fign}[3]{%
  \begin{figure}[tbp]\centering\includegraphics{#1}\caption{#2}\label{#3}\end{figure}}
